\documentclass{beamer}

\usepackage{graphicx}
\usepackage{textpos}
\usepackage{listings}
\usepackage{lstautogobble}

\usetheme{Madrid}
\useoutertheme{miniframes} % Alternatively: miniframes, infolines, split

\newcommand{\xdownarrow}[1]{%
  {\left\downarrow\vbox to #1{}\right.\kern-\nulldelimiterspace}
}

% Setup the university's color pallette
\definecolor{UIUCorange}{RGB}{19, 41, 75} % UBC Blue (primary)
\definecolor{UIUCblue}{RGB}{232, 74, 39} % UBC Grey (secondary)

\definecolor{codegreen}{rgb}{0,0.6,0}
\definecolor{codegray}{rgb}{0.5,0.5,0.5}
\definecolor{codepurple}{rgb}{0.58,0,0.82}
\definecolor{backcolour}{rgb}{0.95,0.95,0.92}

\lstdefinestyle{python}{
  backgroundcolor=\color{backcolour},   
  commentstyle=\color{codegreen},
  keywordstyle=\color{magenta},
  numberstyle=\tiny\color{codegray},
  stringstyle=\color{codepurple},
  basicstyle=\ttfamily\footnotesize,
  breakatwhitespace=false,         
  belowskip=-0.5em,
  breaklines=true,                 
  captionpos=b,                    
  keepspaces=true,                 
  numbers=left,                    
  numbersep=5pt,                  
  showspaces=false,                
  showstringspaces=false,
  showtabs=false,                  
  tabsize=2
}

\lstset{style=python}

\AtBeginSection[]{
  \begin{frame}
    \vfill
    \centering
    \begin{beamercolorbox}[sep=8pt,center,shadow=true,rounded=true]{title}
      \usebeamerfont{title}\insertsectionhead\par%
    \end{beamercolorbox}
    \vfill
  \end{frame}
}
% Setup the university's color pallette
\definecolor{UIUCorange}{RGB}{19, 41, 75} % UBC Blue (primary)
\definecolor{UIUCblue}{RGB}{232, 74, 39} % UBC Grey (secondary)


\setbeamercolor{palette primary}{bg=UIUCorange,fg=white}
\setbeamercolor{palette secondary}{bg=UIUCblue,fg=white}
\setbeamercolor{palette tertiary}{bg=UIUCblue,fg=white}
\setbeamercolor{palette quaternary}{bg=UIUCblue,fg=white}
\setbeamercolor{structure}{fg=UIUCorange} % itemize, enumerate, etc
\setbeamercolor{section in toc}{fg=UIUCblue} % TOC sections

\setbeamercolor{subsection in head/foot}{bg=UIUCorange,fg=UIUCblue}
\setbeamercolor{subsection in head/foot}{bg=UIUCorange,fg=UIUCblue}

\usepackage[utf8]{inputenc}
\usepackage{graphicx}

%Information to be included in the title page:
\title{\textbf{Topic 2: Vars and Expressions}}
\author{\textbf{David H Smith IV}}
\institute[\textbf{UIUC}]{\textbf{University of Illinois Urbana-Champaign}}
\date{\textbf{Thur, Aug 26 2021}}

\setbeamertemplate{title page}[default][colsep=-4bp,rounded=true]
\addtobeamertemplate{title page}{\vspace{3\baselineskip}}{}
\addtobeamertemplate{title page}{
  \begin{textblock*}{\paperwidth}(-1.0em, -1.2em)
    \includegraphics[width=\paperwidth, height=\paperheight]{imgs/uiuc.png}
  \end{textblock*} 
}{}

\begin{document}

\frame{\titlepage}

\section{Announcements}

\begin{frame}
  \begin{enumerate}
    \item Homework 2 (Part 2) and Post-Reading for Topic 3 (Part 1) are posted and due Friday.
    \item Participation Topic 3 (Part 1) is due Friday
    \item Challenge Activities for Topic 2 are due Sunday.
  \end{enumerate}
\end{frame}

\section{Review Poll Questions}

%
%
%
\begin{frame}[fragile]
  \frametitle{Poll Question: }
  How many of the following characters are \textbf{visible} on the screen?\\
  \begin{lstlisting}[language=Python, autogobble]
  print("\t\\n\\\t")
  \end{lstlisting}
  \vfill
  \begin{minipage}{.48\textwidth}
    \begin{enumerate}[A]
      \item 1
      \item 2
      \item 3
      \item 4
      \item 5
    \end{enumerate}
  \end{minipage}
  \begin{minipage}{.48\textwidth}
  \end{minipage}
\end{frame}


\section{Math Operators}

%
%
%
\begin{frame}[fragile]
  \frametitle{Poll Question: Multiplication}
  \vfill
  What is the result of the following?
  \begin{lstlisting}[language=Python, autogobble]
  x = 2(10 + 2)
  print(x)
  \end{lstlisting}
  \vfill
  \begin{enumerate}[A]
    \item 24
    \item SyntaxError
    \item TypeError
    \item ValueError
  \end{enumerate}
\end{frame}

%
%
%
\begin{frame}[fragile]
  \frametitle{Poll Question: Math}
  \vfill
  What is the value of \lstinline|y| after this code executes?
  \begin{lstlisting}[language=Python, autogobble]
  x = 2
  y = x + 3
  x = 3
  \end{lstlisting}
  \vfill
  \begin{enumerate}[A]
    \item 2 
    \item 3
    \item 5
    \item 8
    \item 10
  \end{enumerate}
\end{frame}


%
%
%
\begin{frame}[fragile]
  \frametitle{Poll Question: Math}
  \vfill
  What is the value of \lstinline|y| after this code executes?
  \begin{lstlisting}[language=Python, autogobble]
  x = 7
  y = x
  x = x + 2
  \end{lstlisting}
  \vfill
  \begin{enumerate}
    \item 2
    \item 5 
    \item 7
    \item 9
  \end{enumerate}
\end{frame}

%
% Key point of this slide: Don't get thrown off by the syntactic sugar of Python
%
\begin{frame}[fragile]
  \frametitle{Poll Question: Math}
  \vfill
  What is the value of this expression?
  \begin{lstlisting}[language=Python, autogobble]
  -3 ** 2
  \end{lstlisting}
  \vfill
  \begin{enumerate}[A]
    \item -9
    \item -8
    \item 8
    \item 9
  \end{enumerate}
\end{frame}

%
%
%
\begin{frame}
  \frametitle{Order of Operations in Python}
  \centering
  \begin{minipage}{0.49\textwidth}
    \begin{itemize}
      \item Parentheses 
      \item Exponentiation
      \item Positive and negative
      \item Multiplication, Division, Modulo
      \item Addition, Subtraction 
    \end{itemize}
  \end{minipage}
  \begin{minipage}{0.2\textwidth}
    \centering
    Highest\\
    $\xdownarrow{2cm}$\\
    Lowest\\
  \end{minipage}
  \vfill
  \textbf{Note:} Python evaluates from left to right within a precedence level
\end{frame}


%
%
%
\begin{frame}[fragile]
  \frametitle{Poll Question: More Math Operators}
  \vfill
  Which computes how many (whole) apples I can give to each friend?
  \begin{enumerate}
    \item \lstinline|num_apples / num_friends|
    \item \lstinline|num_friends // num_apples|
    \item \lstinline|num_apples // num_friends|
    \item \lstinline|num_friends % num_appples|
    \item \lstinline|num_apples % num_friends|
  \end{enumerate}
\end{frame}

%
%
%
\begin{frame}[fragile]
  \frametitle{Poll Question: More Math Operators}
  \vfill
  Which computes how many (whole) apples you have left over if you give \lstinline|num_apples| to \lstinline|num_friends|?
  \begin{enumerate}
    \item \lstinline|num_apples / num_friends|
    \item \lstinline|num_friends // num_apples|
    \item \lstinline|num_apples // num_friends|
    \item \lstinline|num_friends % num_appples|
    \item \lstinline|num_apples % num_friends|
  \end{enumerate}
\end{frame}


%
%
%
\begin{frame}[fragile]
  \frametitle{Division, Floor Division, and Modulo}
  \begin{enumerate}
    \item Division operator (/) gives best approximation to true result and \textit{always return a float}.
    \item Floor division (//) rounds down the closet whole number. The type of the result will follow the normal rules.
    \item Modulo operator(\%) performs a division and returns the remainder. The type of the result will always be the same.
    \item For any numbers x and y, the following equality holds: \lstinline|(y == (y // x) * x + (y % x))|
  \end{enumerate}
\end{frame}


\section{Math Module}

%
%
%
\begin{frame}[fragile]
  \frametitle{Poll Question: More Math Operators}
  \vfill
  Which of the following will print the value of $\pi$?
  \begin{enumerate}[A]
    \item \lstinline|print(math.pi)|
    \item \lstinline|print(pi)|
    \item import math.pi\\print(math.pi)
    \item import math\\print(math.pi)
  \end{enumerate}
\end{frame}


%
%
%
\begin{frame}[fragile]
  \frametitle{Poll Question: Rounding}
  What is the result of this code if the user types in \lstinline|4.51| and \lstinline|5.9|?
  \begin{lstlisting}[language=Python, autogobble]
x = math.ceil(float(input()))
y = math.floor(float(input()))
print(x + y)
  \end{lstlisting}
  \vfill
  \begin{enumerate}[A]
    \item SyntaxError
    \item NameError
    \item 10
    \item 10.0
  \end{enumerate}
\end{frame}

%
%
%
\begin{frame}[fragile]
  \frametitle{Poll Question: Rounding}
  What is the result of this code if the user types in \lstinline|4.1| and \lstinline|5.9|?
  \begin{lstlisting}[language=Python, autogobble]
import math
x = math.ceil(float(input()))
y = math.floor(float(input()))
print(x + y)
  \end{lstlisting}
  \vfill
  \begin{enumerate}[A]
    \item 10
    \item 10.0
  \end{enumerate}
\end{frame}

\begin{frame}[fragile]
  \frametitle{Math Module}
  \begin{enumerate}[A]
    \item Should I memorize (operators, function, modules, module functions, etc.)?
    \pause
    \item Yes and no. 
    \pause
    \item Google + \lstinline|help()| function are your friends
    \pause
    \item Modules vs Scripts: Modules are just (for the purposes of what we've discussed so far) scripts that someone else wrote that you can use in your own scripts.
    \pause
  \end{enumerate}
  \vfill
  \pause
  To get information on a module:
  \begin{lstlisting}[language=Python, autogobble]
  import math
  help(math)
  \end{lstlisting}
\end{frame}

\begin{frame}
  \frametitle{\lstinline|__name__| and \lstinline|"__main__"|}

  \begin{enumerate}
    \item When you run a script in python it gets a few "environment variates".
    \item For the script you run (e.g, \lstinline|test.py|) the \lstinline|__name__| variable will always be \lstinline|"__main__"|.
    \item For any scripts/modules you import \lstinline|__name__| variable will always be the name of that script/module.
  \end{enumerate}
\end{frame}

\section{Lab}


\section{Lab 1}

\begin{frame}
  \frametitle{Python + TextEditor}
  \begin{enumerate}
    \item Everyone should be able to open a zsh terminal (Mac), powershell cmd (Windows), or Git Bash (Either) and type: \lstinline|python|.
    \item You should be greeted by a Python command line (e.g., \lstinline|>>>|).
  \end{enumerate}
\end{frame}

\begin{frame}
  \frametitle{Git Setup}
  \begin{enumerate}
    \item \lstinline|git config|: Sets up your local 
    \item All of these local setting will be represented in your local "commit history".
    \item When you push that commit history to GitHub (a remote repository) they will also be reflected there.
  \end{enumerate}
\end{frame}

\begin{frame}
  \frametitle{GitHub Setup}
  \begin{enumerate}
    \item \textbf{Secure Socket Shell (SSH) Key: } A "password" for remote access to a repository.
    \item Follow the instructions listed on the link in the lab instructions.
    \item The key generation command it will give you a path to a file that ends in \lstinline|.pub|
    \item Open the aformentioned file and copy-paste it's contents into the key section of GitHub.
  \end{enumerate}
\end{frame}

\begin{frame}
  \frametitle{Cloning A Repository}
\end{frame}

%
%
%
\begin{frame}
  \begin{enumerate}
    \item Homework 2 (Part 2) and Post-Reading for Topic 3 (Part 1) are posted and due Friday.
    \item Participation Topic 3 (Part 1) is due Friday
    \item Challenge Activities for Topic 2 are due Sunday.
    \item Lab 1 (100pts) will be due Sunday, Sept 4th:
      \begin{itemize}
        \item Setup (50pts).
        \item README.md with some stuff filled in (25pts).
        \item \lstinlne|hello-world.py| passes the test cases (25pts).
      \end{itemize}
  \end{enumerate}
\end{frame}


\end{document}
